\section{Giới thiệu}\label{sec:intro}
Dịch vụ bản đồ giúp tìm đường và địa điểm trong lúc di chuyển. Nếu nhiều chuyến thường kết thúc gần một bệnh viện, bên quan sát có thể phỏng đoán về sức khỏe hoặc thói quen dù người dùng chỉ tới thăm người thân. Tính đặc trưng của dấu vết di chuyển làm rủi ro này đáng quan tâm~\cite{mobilityunique}.

Câu hỏi chúng tôi đặt ra là sau cả chuyến đi, chuỗi trao đổi ấy giúp nhà
cung cấp phân biệt hai quỹ đạo thật đến mức nào? Geo-indistinguishability
(Geo-I) cho ta điểm xuất phát bằng cách giới hạn tỷ số xác suất quan sát cùng
một đầu ra tại hai vị trí thật~\cite{geoi}. Để mở rộng lập luận ra cả phiên,
ta phải tính cả những lần thiết bị đọc lại GPS để quyết định giữ hay đổi điểm
tham chiếu. Chi phí của quan sát mới cũng là cơ sở của các cơ chế bảo vệ dùng
dự đoán~\cite{predictive}.

Về phía người dùng, điều quan trọng là danh sách địa điểm (POI) nào thực sự
hiện trên màn hình. Truy vấn từ nhiều vị trí có thể thu về đủ địa điểm để
điện thoại xếp hạng theo nhu cầu thật~\cite{ev}. Vì vậy, chúng tôi so danh sách
người dùng nhận được với đáp án cho cùng yêu cầu tại GPS thật trên toàn bộ
danh mục. Cách đo này áp dụng được khi đầu ra công khai là một tọa độ, một
tập điểm thật và giả, hoặc một chuỗi có thêm truy vấn giả. Các công trình về
phát hành quỹ đạo ngoại tuyến có mục tiêu sử dụng khác~\cite{transprotect,semantic,fake,sok}.

Chúng tôi xác định mười kịch bản theo thông tin bên quan sát muốn suy ra
và phần chuỗi trao đổi (transcript) họ được thấy.

\smallskip\noindent\textbf{Vị trí và đường đã đi.}
Từ một gói, bên quan sát ước lượng GPS hiện tại (S1). Với chuỗi dài hơn,
họ có thể tìm nơi dừng (S2) hoặc khôi phục vị trí ở từng mốc thời gian (S3).
Ở S2, dữ liệu chỉ cho phép dùng tọa độ trung bình trong phút cuối làm đại
diện cho nơi dừng. Khi cho xem toàn phiên, ta còn đo được khả năng suy ra
điểm xuất phát (S9) và điểm kết thúc (S10). Đây là năm tác vụ định vị chính.

\smallskip\noindent\textbf{Liên kết và dự đoán.}
Hai phiên có thuộc về cùng một người hay phương tiện (S4)? Dữ liệu chỉ có
nhãn nhóm tuyến, nên phép thử là xác định hai phiên có cùng nhóm hay không.
Với suy luận tương lai, chúng tôi xét hai quyết định khác nhau trên mạng
đường. S5 dự đoán nhánh kế tiếp trong hai nhánh khả dĩ. S6 chọn một trong
hai đích đến khi đã thấy người dùng rẽ.

\smallskip\noindent\textbf{Nội dung và người đi cùng.}
Ở S7, ta giữ nguyên chuyến đi rồi đổi yêu cầu trên thiết bị để kiểm tra xem
nội dung yêu cầu có làm thay đổi transcript hay không. S8 xét một nguồn
thông tin khác mà nhà cung cấp có thể có. Khi căn hai luồng truy vấn theo
thời gian, quan sát thêm người đi cùng giúp định vị người dùng đến mức nào?


GeoI-BQ giữ điểm tham chiếu đã bảo vệ trên mạng đường. Từ lịch sử các điểm
tham chiếu, thiết bị ước lượng một phân phối vị trí (belief) để chọn cả đoạn
truy vấn kế tiếp. Nghiên cứu có ba đóng góp.
\begin{itemize}
\item Một khung mô tả từng kịch bản theo mục tiêu suy luận, lịch sử được quan sát và dữ liệu tham chiếu để chấm.
\item Một cơ chế truy vấn theo đoạn dựa trên Geo-I, với cận tiết lộ cả phiên và điều kiện bảo đảm chất lượng dịch vụ.
\item Đối chiếu sáu phương án trên cùng tác vụ, có khoảng tin cậy, chi phí và các thước đo gốc áp dụng được.
\end{itemize}
Trong thử nghiệm, thiết bị đã có toàn bộ danh mục tĩnh và phương án nền
công khai đạt Recall nhỉnh hơn cơ chế thích nghi. Hai điều kiện này ảnh hưởng
đến cách diễn giải lợi ích của truy vấn trực tuyến.
Các phần~\ref{sec:related}--\ref{sec:conclusion} trình bày nghiên cứu liên quan,
phương pháp và chứng minh, kết quả thực nghiệm và kết luận.
